  \recommended \item
    已知 \( a,b,c\in\Re \)，要取什么样的 \( d \) 才能使这个矩阵
    的秩为一？
    \begin{equation*}
      \begin{mat}
         a  &b  \\
         c  &d
      \end{mat}
    \end{equation*}
    \begin{answer}
      若 \( a\neq 0 \)，则取 \( d=(c/a)b \) 可使第二行是第一行的倍数，
      具体地说，是第一行的 \( c/a \) 倍。
      若 \( a=0 \) 且 \( b=0 \)，则 \( d \) 取任意非 \( 0 \) 值都能保证
      第二行非零。
      若 \( a=0 \) 且 \( b\neq 0 \) 且 \( c=0 \)，则 \( d \) 取任何值都行，
      因为该矩阵自动具有秩一（即使取 $d=0$）。
      最后，若 \( a=0 \) 且 \( b\neq 0 \) 且 \( c\neq 0 \)，则 \( d \) 取
      任何值都不够，因为该矩阵必定有秩二。
    \end{answer}
  \item
    求这个矩阵的列秩。
    \begin{equation*}
      \begin{mat}[r]
        1  &3  &-1  &5  &0  &4  \\
        2  &0  &1   &0  &4  &1
      \end{mat}
    \end{equation*}
    \begin{answer}
      列秩是二。
      一种看法是直接观察\Dash 列空间由高度为二的列组成，因而其维数
      至少可以是二，而且我们容易找到两列，它们合起来组成一个
      线性无关的集合（例如第四列与第五列）。
      另一种看法是回忆列秩等于行秩，然后施行高斯消元法，
      结果剩下两个非零行。
    \end{answer}
  \item
    证明：一个至少有一个解的线性方程组至多有一个解当且仅当
    其系数矩阵的秩等于其列数。
    \begin{answer}
      我们应用\nearbytheorem{th:RankVsSoltnSp}。
      线性方程组的系数矩阵 $A$ 的列数等于未知数的个数 $n$。
      一个至少有一个解的线性方程组至多有一个解当且仅当相应的
      齐次方程组的解空间维数为零（回忆：~在
      “$\text{通解}=\text{特解}+\text{齐次解}$”这一方程
       $\vec{v}=\vec{p}+\vec{h}$ 中，只要这样的 $\vec{p}$ 存在，
       解 $\vec{v}$ 唯一当且仅当向量 $\vec{h}$ 唯一，即 $\vec{h}=\zero$）。
       而由该定理，这意味着 $n=r$。
    \end{answer}
  \recommended \item
    如果一个矩阵是 \( \nbym{5}{9} \) 的，它的行集合与列集合当中，
    哪一个必定是线性相关的？
    \begin{answer}
      列的集合必定是线性相关的，因为该矩阵的秩至多是五，而列却有
      九个。
    \end{answer}
  \item
    举一个例子，说明尽管两者维数相同，一个矩阵的行空间与列空间
    却不必相等。
    它们有可能相等吗？
    \begin{answer}
      它们几乎没有相等的可能，因为行空间是行向量的集合，而列空间是
      列向量的集合（除非该矩阵是 \( \nbyn{1} \) 的，此时两个空间
      必定相等）。

      \textit{注记。}
      考察
      \begin{equation*}
        A=\begin{mat}[r]
            1  &3  \\
            2  &6
          \end{mat}
      \end{equation*}
      并注意行空间是 \( \rowvec{1 &3} \) 的所有倍数组成的集合，
      而列空间由以下列向量的倍数组成
      \begin{equation*}
        \colvec[r]{1 \\ 2}
      \end{equation*}
      所以我们也不能论证说这两个空间必定恰好互为转置。
    \end{answer}
  \item
    证明集合
    \( \set{(1,-1,2,-3),(1,1,2,0),(3,-1,6,-6)} \) 与
    \( \set{(1,0,1,0),(0,2,0,3)} \) 的张成不相同。
    顺便问一句，这里的向量空间是什么？
    \begin{answer}
      首先，该向量空间是在自然运算下的实数四元组的集合。
      尽管这不是四宽的行向量的集合，但差别不大\Dash 它与那个集合
      是“相同的”。
      所以我们把四元组当作四宽的向量来处理。

      有了这一点，要看出 \( (1,0,1,0) \) 不在第一个集合的张成中，
      一种办法是注意这个化简
      \begin{equation*}
        \begin{mat}[r]
          1  &-1  &2  &-3  \\
          1  &1   &2  &0   \\
          3  &-1  &6  &-6
        \end{mat}
        \grstep[-3\rho_1+\rho_3]{-\rho_1+\rho_2}
        \repeatedgrstep{-\rho_2+\rho_3}
        \begin{mat}[r]
          1  &-1  &2  &-3  \\
          0  &2   &0  &3   \\
          0  &0   &0  &0
        \end{mat}
      \end{equation*}
      以及这个化简
      \begin{equation*}
        \begin{mat}[r]
          1  &-1  &2  &-3  \\
          1  &1   &2  &0   \\
          3  &-1  &6  &-6  \\
          1  &0   &1  &0
        \end{mat}
        \grstep[-3\rho_1+\rho_3 \\ -\rho_1+\rho_4]{-\rho_1+\rho_2}
        \repeatedgrstep[-(1/2)\rho_2+\rho_4]{-\rho_2+\rho_3}
        \repeatedgrstep{\rho_3\leftrightarrow\rho_4}
        \begin{mat}[r]
          1  &-1  &2  &-3  \\
          0  &2   &0  &3   \\
          0  &0   &-1 &3/2 \\
          0  &0   &0  &0
        \end{mat}
      \end{equation*}
      得到的矩阵在秩上不同。
      这意味着把 $(1,0,1,0)$ 添加到前三个四元组的集合中会增大该
      集合的秩，从而增大其张成。
      因此 $(1,0,1,0)$ 原本就不在这个张成中。
    \end{answer}
  \recommended \item
    证明这个列向量的集合
    \begin{equation*}
      \set{\colvec{d_1 \\ d_2 \\ d_3}
           \suchthat
           \text{存在 $x$、$y$ 与 $z$ 使得：\ }
           \begin{linsys}{3}
             3x  &+  &2y  &+  &4z  &=   &d_1   \\
              x  &   &    &-  &z   &=   &d_2   \\
             2x  &+  &2y  &+  &5z  &=   &d_3
           \end{linsys}
           }
    \end{equation*}
    是 \( \Re^3 \) 的一个子空间。
    求一个基。
    \begin{answer}
      它是子空间，因为它是系数矩阵
      \begin{equation*}
        \begin{mat}[r]
          3  &2  &4  \\
          1  &0  &-1 \\
          2  &2  &5
        \end{mat}
      \end{equation*}
      的列空间。
      为了求列空间的一个基，
      \begin{equation*}
        \set{c_1\colvec[r]{3 \\ 1 \\ 2}
             +c_2\colvec[r]{2 \\ 0 \\ 2}
             +c_3\colvec[r]{4 \\ -1 \\ 5}
             \suchthat c_1,c_2,c_3\in\Re}
      \end{equation*}
      我们通过转置、化简
      \begin{equation*}
         \begin{mat}[r]
           3  &1  &2  \\
           2  &0  &2  \\
           4  &-1 &5
         \end{mat}
         \grstep[-(4/3)\rho_1+\rho_3]{-(2/3)\rho_1+\rho_2}
         \repeatedgrstep{-(7/2)\rho_2+\rho_3}
         \begin{mat}[r]
           3  &1     &2  \\
           0  &-2/3  &2/3  \\
           0  &0     &0
         \end{mat}
      \end{equation*}
      并略去零行，再转置回去，来消去张成集中三个列向量之间的
      线性关系。
      \begin{equation*}
        \sequence{ \colvec[r]{3 \\ 1 \\ 2},
                   \colvec[r]{0 \\ -2/3 \\ 2/3}  }
      \end{equation*}
    \end{answer}
  \item
    证明转置运算是
    线性的：\index{transpose!is linear}
    \begin{equation*}
      \trans{(rA+sB)}  = r\trans{A}+s\trans{B}
    \end{equation*}
    其中 \( r,s\in\Re \) 且 \( A,B\in\matspace_{\nbym{m}{n}} \)。
    \begin{answer}
      我们可以通过直接的计算来完成。
      \begin{align*}
        \trans{(rA+sB)}
        &=\trans{\begin{mat}
             ra_{1,1}+sb_{1,1}  &\ldots &ra_{1,n}+sb_{1,n} \\
                          &\vdots            \\
             ra_{m,1}+sb_{m,1}  &\ldots &ra_{m,n}+sb_{m,n}
                \end{mat}  }  \\
        &=\begin{mat}
           ra_{1,1}+sb_{1,1}  &\ldots &ra_{m,1}+sb_{m,1}  \\
                            &\vdots                           \\
           ra_{1,n}+sb_{1,n}  &\ldots &ra_{m,n}+sb_{m,n}
         \end{mat}           \\
        &=\begin{mat}
           ra_{1,1}  &\ldots &ra_{m,1}  \\
                    &\vdots                       \\
           ra_{1,n}  &\ldots &ra_{m,n}
         \end{mat}
        +\begin{mat}
           sb_{1,1}  &\ldots &sb_{m,1}  \\
                    &\vdots                       \\
           sb_{1,n}  &\ldots &sb_{m,n}
         \end{mat}           \\
        &=r\trans{A}+s\trans{B}
      \end{align*}
    \end{answer}
  \recommended \item
    在本小节中我们已经证明高斯化简能求出行空间的一个基。
    \begin{exparts}
      \partsitem 证明这个基不是唯一的\Dash 不同的化简可能给出
        不同的基。
      \partsitem 构造行空间相等但行数不相等的矩阵。
      \partsitem 证明两个矩阵有相同的行空间当且仅当经过
        高斯–若尔当消元后它们有相同的非零行。
    \end{exparts}
    \begin{answer}
      \begin{exparts}
        \partsitem 这些化简给出不同的基。
          \begin{equation*}
            \begin{mat}[r]
              1  &2  &0  \\
              1  &2  &1
            \end{mat}
            \grstep{-\rho_1+\rho_2}
            \begin{mat}[r]
              1  &2  &0  \\
              0  &0  &1
            \end{mat}
            \qquad
            \begin{mat}[r]
              1  &2  &0  \\
              1  &2  &1
            \end{mat}
            \grstep{-\rho_1+\rho_2}
            \repeatedgrstep{2\rho_2}
            \begin{mat}[r]
              1  &2  &0  \\
              0  &0  &2
            \end{mat}
          \end{equation*}
        \partsitem 一个简单的例子是这个。
          \begin{equation*}
            \begin{mat}[r]
              1  &2  &1  \\
              3  &1  &4
            \end{mat}
            \qquad
            \begin{mat}[r]
              1  &2  &1  \\
              3  &1  &4  \\
              0  &0  &0
            \end{mat}
          \end{equation*}
          这是一个稍微不那么简单的例子。
          \begin{equation*}
            \begin{mat}[r]
              1  &2  &1  \\
              3  &1  &4
            \end{mat}
            \qquad
            \begin{mat}[r]
              1  &2  &1  \\
              3  &1  &4  \\
              2  &4  &2  \\
              4  &3  &5
            \end{mat}
          \end{equation*}
        \partsitem
          % We give two proofs.
          % The first is a bit lower-level but more constructive.
          % The second is higher level.

          % \textit{First proof.}
          % Assume that \( A \) and \( B \) are matrices with equal row spaces.
          % Construct a matrix \( C \) with the rows of \( A \) above the rows
          % of \( B \), and another matrix \( D \) with the rows of \( B \)
          % above the rows of \( A \).
          % \begin{equation*}
          %   C=\begin{mat}
          %        A \\ B
          %     \end{mat}
          %   \qquad
          %   D=\begin{mat}
          %        B \\ A
          %     \end{mat}
          % \end{equation*}
          % Observe that \( C \) and \( D \) are row-equivalent (via a sequence
          % of row-swaps) and so Gauss-Jordan reduce to the same reduced
          % echelon form matrix.

          % Because the row spaces are equal, the rows of \( B \)
          % are linear combinations of the rows of \( A \) so Gauss-Jordan
          % reduction on \( C \) simply turns the rows of \( B \) to zero rows
          % and thus the nonzero rows of $C$ are just the nonzero rows obtained
          % by Gauss-Jordan reducing \( A \).
          % The same can be said for the matrix \( D \)\Dash Gauss-Jordan
          % reduction on \( D \) gives the same non-zero rows as are produced
          % by reduction on \( B \) alone.
          % Therefore,
          % \( A \) yields the same nonzero rows as \( C  \), which yields
          % the same nonzero rows as \( D \), which yields the same nonzero
          % rows as \( B \).

          % \textit{Second proof (N~Schenker).}
          由于 \( A \) 与 \( B \) 的行空间相等，所以二者行等价。
          由于每个行等价类都包含唯一的简化阶梯形
          （定理 One.III.\ref{th:ReducedEchelonFormIsUnique}），
          \( A \) 的简化阶梯形必定等于 \( B \) 的简化阶梯形。
      \end{exparts}
    \end{answer}
  \item
    对于 \nearbyremark{rem:MunkresCommment}，在 \( r \) 大于 \( n \) 的
    情形下为什么不成问题？
    \begin{answer}
      它不可能更大。
    \end{answer}
  \item
    证明 \( \nbym{m}{n} \) 矩阵的行秩至多是 \( m \)。
    有更好的界吗？
    \begin{answer}
      极大线性无关集合中行的个数不会超过行的个数。
      一个更好的界（一般情形下这是最佳可能的界）是 \( m \) 与 \( n \)
      中的最小者，因为行秩等于列秩。
     \end{answer}
  \item
    证明一个矩阵的秩等于其转置的秩。
    \begin{answer}
      由于矩阵 \( A \) 的行是 \( \trans{A} \) 的列，所以 \( A \) 的
      行空间的维数等于 \( \trans{A} \) 的列空间的维数。
      而 \( A \) 的行空间的维数就是 \( A \) 的秩，并且 \( \trans{A} \) 的
      列空间的维数就是 \( \trans{A} \) 的秩。
      于是这两个秩相等。
    \end{answer}
  \item
    判断真假：~一个矩阵的列空间等于其转置的行空间。
    \begin{answer}
      假的。
      前者是列的集合，而后者是行的集合。

      然而这个例子
      \begin{equation*}
         A=\begin{mat}[r]
             1  &2  &3  \\
             4  &5  &6
           \end{mat},
         \qquad
         \trans{A}=\begin{mat}[r]
             1  &4  \\
             2  &5  \\
             3  &6
           \end{mat}
      \end{equation*}
      表明，一旦我们对“相同”有了形式化的含义，就可以在这里应用它：
      \begin{equation*}
        \colspace{A}=\spanof{\set{
                        \colvec[r]{1 \\ 4},
                        \colvec[r]{2 \\ 5},
                        \colvec[r]{3 \\ 6} }  }
      \end{equation*}
      而
      \begin{equation*}
         \rowspace{\trans{A}}=\spanof{\set{
                                 \rowvec{1 &4},
                                 \rowvec{2 &5},
                                 \rowvec{3 &6} }  }
      \end{equation*}
      彼此是“相同的”。
    \end{answer}
  \recommended \item
    我们已经看到行变换可能改变列空间。
    它必定会吗？
    \begin{answer}
      不会。
      这里，高斯消元法没有改变列空间。
      \begin{equation*}
        \begin{mat}[r]
          1  &0  \\
          3  &1
        \end{mat}
        \grstep{-3\rho_1+\rho_2}
        \begin{mat}[r]
          1  &0  \\
          0  &1
        \end{mat}
      \end{equation*}
    \end{answer}
  \item
    证明一个线性方程组有解当且仅当该方程组的系数矩阵与
    其增广矩阵有相同的秩。
    \begin{answer}
      线性方程组
      \begin{equation*}
        c_1\vec{a}_1+\dots+c_n\vec{a}_n=\vec{d}
      \end{equation*}
      有解当且仅当 \( \vec{d} \) 在集合 \( \set{\vec{a}_1,\dots,\vec{a}_n} \)
      的张成中。
      这成立当且仅当增广矩阵的列秩等于系数矩阵的列秩。
      由于秩等于列秩，方程组有解当且仅当其增广矩阵的秩等于其
      系数矩阵的秩。
    \end{answer}
  \item
    一个 \( \nbym{m}{n} \) 矩阵具有
    \definend{行满秩}\index{full row rank}\index{row rank!full}，
    若其行秩
    为 \( m \)；它具有
    \definend{列满秩}\index{full column rank}\index{column!rank!full}，
    若其列秩为 \( n \)。
    \begin{exparts}
      \partsitem 证明一个矩阵能同时有行满秩与列满秩仅当它是方阵。
      \partsitem 证明：系数矩阵为 \( A \) 的线性方程组对右侧任意的
        \( d_1 \)、\ldots、\( d_n \) 都有解当且仅当 \( A \) 行满秩。
      \partsitem 证明齐次方程组有唯一解当且仅当其系数矩阵
        $A$ 列满秩。
      \partsitem 证明命题
       “若系数矩阵为 \( A \) 的方程组有任何解则它有唯一解”成立
        当且仅当 \( A \) 列满秩。
    \end{exparts}
    \begin{answer}
      \begin{exparts}
        \partsitem 行秩等于列秩，所以二者各自至多是行数与列数的
          最小者。
          因此只有当行数等于列数时二者才能都满。
          （当然，逆命题不成立：~方阵不必行满秩或列满秩。）
        \partsitem 若 \( A \) 行满秩，则无论右侧是什么，对增广矩阵施行
          高斯消元法结束时每一行都有一个首主元一，并且这些首主元一
          没有一个位于最右边的列（“增广”列）。
          回代随即给出一个解。

          另一方面，如果该线性方程组对某个右侧没有解，只可能是因为
          高斯消元法留下某行，使得它在“增广”竖线左边全为零，
          而在右边有非零元。
          于是，如果 $A$ 对某些右侧没有解，那么 \( A \) 不具有
          行满秩，因为它的某些行已被消去。
        \partsitem 矩阵 \( A \) 列满秩当且
          仅当其列
          组成一个线性无关的集合。
          这等价于诸列之间只存在平凡的线性关系，于是方程组唯一的解是
          每个变量都取 $0$ 的情形。
        \partsitem 矩阵 \( A \) 列满秩当且
          仅当其列的集合是线性无关的，从而构成其张成的一个基。
          这等价于该张成中的所有向量都有唯一的线性表示。
          这就证明了它，
          因为张成中一个向量的任何线性表示都是
          该线性方程组的解。
      \end{exparts}
    \end{answer}
  \item
    如果把高斯消元法改成允许用零去乘某一行，
    \nearbylemma{le:RowSpUnchByGR} 的结论会如何改变？
    \begin{answer}
      行空间不再相同，$B$ 的行空间将是 $A$ 的行空间的一个子空间
      （有可能相等）。
    \end{answer}
  \item
    \( \rank(A) \) 与 \( \rank(-A) \) 之间有什么关系？
    \( \rank(A) \) 与 \( \rank(kA) \) 之间呢？
    \( \rank(A) \)、\( \rank(B) \) 与 \( \rank(A+B) \) 之间，
    如果有关系的话，是什么关系？
    \begin{answer}
      显然 \( \rank(A)=\rank(-A) \)，因为高斯消元法允许我们把一个
      矩阵的所有行乘以 \( -1 \)。
      同样地，当 \( k\neq 0 \) 时我们有 \( \rank(A)=\rank(kA) \)。

      加法更有意思。
      和的秩可以比各加数的秩小。
      \begin{equation*}
        \begin{mat}[r]
          1  &2  \\
          3  &4
        \end{mat}
        +
        \begin{mat}[r]
         -1  &-2  \\
         -3  &-4
        \end{mat}
        =
        \begin{mat}[r]
         0   &0   \\
         0   &0
        \end{mat}
      \end{equation*}
      和的秩也可以比各加数的秩大。
      \begin{equation*}
        \begin{mat}[r]
          1  &2  \\
          0  &0
        \end{mat}
        +
        \begin{mat}[r]
          0  &0   \\
          3  &4
        \end{mat}
        =
        \begin{mat}[r]
         1   &2   \\
         3   &4
        \end{mat}
      \end{equation*}
      但存在一个上界（除矩阵尺寸之外的上界）。
      一般地，\( \rank(A+B)\leq\rank(A)+\rank(B) \)。

      为了证明它，注意我们可以用两种方式之一对 \( A+B \) 实施
      高斯消去法：
      可以先把 \( A \) 加到 \( B \) 上，然后施行适当的化简步骤序列
      \begin{equation*}
        (A+B)
        \grstep{\text{步骤}_1}
        \cdots
        \grstep{\text{步骤}_k}
        \text{阶梯形}
      \end{equation*}
      或者也可以分别在 \( A \) 与 \( B \) 上实施 \( \text{步骤}_1 \) 到
      \( \text{步骤}_k \)，然后再相加，得到同样的结果。
      在第二种方式中我们最终能得到的最大秩显然是两个秩之和。
      （上面的矩阵给出了两种可能性，
       $\rank(A+B)<\rank(A)+\rank(B)$ 与 $\rank(A+B)=\rank(A)+\rank(B)$，
       都发生的例子。）
    \end{answer}
    % Relationship between rank(A), rank(B) and rank(cA+dB)?
\end{exercises}























